\section{Discussion and Future Work}\label{sec:discussion}

\paragraph{Formal soundness.}
The most important open piece is a small-step operational semantics
with a soundness theorem: if a program never uses a poisoned
reference, never has two overlapping mutable arguments at a call,
and never returns a reference that does not trace back to a
parameter, then its reference operations preserve type and memory
safety. The rules in Sect.~\ref{sec:semantics} are designed to make
this provable.

\paragraph{Opaque effects.}
The interface for native functions, closure captures, and
dynamic-dispatch targets is small but part of the trusted base. An
opaque operation declares which of its reference arguments each
returned reference is derived from; the runtime treats the operation
as a black-box step that produces a fresh child node under that
argument. This is enough to keep the discipline working at the
boundary, but correctness then depends on the declared model. A
natural follow-up is a unified treatment of all three cases together
with the soundness condition an opaque model must satisfy.

\paragraph{Deployment modes.}
Because invalidation is a one-way bit, the discipline can be deployed
in several modes. The most aggressive is \emph{always-on}, useful as
defense in depth on bytecode whose verifier outcome is not trusted.
A natural relaxation is \emph{trace-and-replay}: record the relevant
events during execution and run the check on the recorded trace
later, in parallel with other work. Production deployments will
likely mix the two depending on the trust they place in the static
verifier on each code path.

\paragraph{Limitations.}
Trusted models can silently break the guarantees if they misdescribe
an opaque effect. Vector indices collapse to a single abstract label,
so the discipline cannot distinguish two references into different
elements of the same vector. And when a violation is detected, the
runtime reports the event but does not yet emit a derivation trace
explaining which earlier operation poisoned the reference; such a
trace would help in debugging.

\section{Related Work}\label{sec:related}

\paragraph{Rust-style static borrow checking.}
The closest static relative is Rust's borrow checker, whose
metatheory was given by RustBelt~\cite{rustbelt}. Both target
stack-allocated state, with no first-class notion of globally-stored
mutable resources, and the only dynamic borrow check in mainstream
Rust is \texttt{RefCell}'s local check on a single cell. Move's
reference-safety rules, in contrast, span locals, globals, and
mutable cross-frame arguments; our discipline enforces them at
execution time.

\paragraph{Per-location dynamic alias models.}
The closest dynamic neighbors are Stacked Borrows~\cite{stackedborrows}
and Tree Borrows~\cite{treeborrows}. Stacked and Tree Borrows model
aliasing at memory locations: each location carries a stack or tree
of borrow tags, and the Miri interpreter checks reads and writes
against this metadata to detect undefined behavior in Rust. Our
model tracks references by their access paths rather than by the
addresses they happen to land on, which matters because Move's
reference rules include events the per-location view does not
naturally express---in particular variant mutation, where a write
to a discriminator invalidates references derived from the prior
variant's projection. We also treat call boundaries as a first-class
construct, with subtree locking and parameter-root transport doing
the work that the static borrow checker does at compile time in
Rust.

\paragraph{Capabilities, separation, runtime verification.}
Capability systems~\cite{capabilities} read naturally onto our model:
a node identity is the capability, and an invalidation is a
revocation. Access paths play the role of a runtime footprint in the
sense of separation logic~\cite{seplogic}, though we do not develop
the logical reading here. The RV literature on monitoring mutable
state~\cite{rvsurvey} provides the broader backdrop; we specialize
to one specific monitor, over the hierarchy of reference derivation.
Move ships with a static reference-safety checker performing abstract
interpretation over a borrow graph, which the runtime discipline is
designed to interoperate with.
